% preamble-exam-zh.tex — 共享 preamble
% 用于所有 exam-zh 模板
%
% 样式策略：
%   屏幕 → 语义彩色（蓝=关键想法，红=易错，绿=路线，灰=提示/教师备注）
%   打印 → 自动降级为黑白（通过 print mode 或 monochrome 选项）

\documentclass{exam-zh}

% ---------- 基础配置 ----------
\examsetup{
  page/size = a4paper,
  page/show-head = false,
  page/show-foot = false,
  sealline/show = false,
  font = times,
  math-font = xits,
}

\usepackage{tcolorbox}
\usepackage{needspace}
\tcbuselibrary{skins,breakable}

% ---------- 打印/屏幕颜色切换 ----------
% 用法：\documentclass[monochrome]{exam-zh} 即可切换为纯黑白
% 注意：不要使用 \if@xxx 放在 makeatother 外部；这里改成普通条件名。
\newif\ifeduprintcolor
\eduprintcolortrue

\makeatletter
\@ifclasswith{exam-zh}{monochrome}{\eduprintcolorfalse}{}
\makeatother

% ---------- 语义颜色定义 ----------
% 屏幕：有色彩区分；打印：降级为灰度
\ifeduprintcolor
  % --- 屏幕彩色模式 ---
  \definecolor{edu-blue}{RGB}{47, 82, 122}       % 关键想法
  \definecolor{edu-blue-bg}{RGB}{243, 246, 252}
  \definecolor{edu-red}{RGB}{180, 40, 40}         % 易错提醒
  \definecolor{edu-red-bg}{RGB}{255, 245, 240}
  \definecolor{edu-green}{RGB}{34, 100, 60}       % 路线图
  \definecolor{edu-green-bg}{RGB}{242, 250, 245}
  \definecolor{edu-orange}{RGB}{160, 90, 0}       % 训练目标
  \definecolor{edu-orange-bg}{RGB}{255, 248, 235}
  \definecolor{edu-gray}{RGB}{100, 100, 100}      % 提示/教师备注
  \definecolor{edu-gray-bg}{RGB}{248, 248, 248}
  \definecolor{edu-purple}{RGB}{90, 50, 120}      % 思考题
  \definecolor{edu-purple-bg}{RGB}{248, 244, 252}
\else
  % --- 打印黑白模式 ---
  \definecolor{edu-blue}{gray}{0.15}
  \definecolor{edu-blue-bg}{gray}{0.96}
  \definecolor{edu-red}{gray}{0.25}
  \definecolor{edu-red-bg}{gray}{0.97}
  \definecolor{edu-green}{gray}{0.20}
  \definecolor{edu-green-bg}{gray}{0.97}
  \definecolor{edu-orange}{gray}{0.25}
  \definecolor{edu-orange-bg}{gray}{0.98}
  \definecolor{edu-gray}{gray}{0.40}
  \definecolor{edu-gray-bg}{gray}{0.97}
  \definecolor{edu-purple}{gray}{0.30}
  \definecolor{edu-purple-bg}{gray}{0.98}
\fi

% ---------- 自定义教学环境 ----------
% 共性：enhanced + breakable（跨页不断裂）

% 原题卡片 — 强边框，突出显示
\newtcolorbox{problemcard}{
  enhanced, breakable,
  colback=edu-blue-bg,
  colframe=edu-blue,
  boxrule=1.2pt,
  arc=0pt,
  left=3mm, right=3mm, top=3mm, bottom=3mm,
}

% 解题路线图 — 左边线 + 浅底
\newtcolorbox{routebox}{
  enhanced, breakable,
  colback=edu-green-bg,
  colframe=edu-green!30,
  boxrule=0pt,
  borderline west={2.5pt}{0pt}{edu-green},
  left=3mm, right=3mm, top=3mm, bottom=3mm,
}

% 关键想法 — 蓝色左边线
\newtcolorbox{keyideabox}{
  enhanced, breakable,
  colback=edu-blue-bg,
  colframe=edu-blue!30,
  boxrule=0pt,
  borderline west={2.5pt}{0pt}{edu-blue},
  left=3mm, right=3mm, top=3mm, bottom=3mm,
}

% 易错提醒 — 红色左边线
\newtcolorbox{mistakebox}{
  enhanced, breakable,
  colback=edu-red-bg,
  colframe=edu-red!20,
  boxrule=0pt,
  borderline west={3pt}{0pt}{edu-red},
  left=3mm, right=3mm, top=2.5mm, bottom=2.5mm,
}

% 提示 — 灰色虚线左边线
\newtcolorbox{hintbox}{
  enhanced, breakable,
  colback=edu-gray-bg,
  colframe=edu-gray!20,
  boxrule=0pt,
  borderline west={2pt}{0pt}{edu-gray, dashed},
  left=3mm, right=3mm, top=2mm, bottom=2mm,
}

% 思考题 — 紫色虚线左边线
\newtcolorbox{thinkbox}{
  enhanced, breakable,
  colback=edu-purple-bg,
  colframe=edu-purple!20,
  boxrule=0pt,
  borderline west={2pt}{0pt}{edu-purple, dashed},
  left=2.5mm, right=2.5mm, top=2.5mm, bottom=2.5mm,
}

% 教师备注 — 灰色左边线，小字
\newtcolorbox{teachernote}{
  enhanced, breakable,
  colback=edu-gray-bg,
  colframe=edu-gray!15,
  boxrule=0pt,
  borderline west={2pt}{0pt}{edu-gray!60},
  left=2.5mm, right=2.5mm, top=2.5mm, bottom=2.5mm,
  fontupper=\small,
  colupper=black!60,
}

% 训练目标 — 橙色左边线，小字
\newtcolorbox{traininggoal}{
  enhanced, breakable,
  colback=edu-orange-bg,
  colframe=edu-orange!15,
  boxrule=0pt,
  borderline west={1.5pt}{0pt}{edu-orange},
  left=2mm, right=2mm, top=1mm, bottom=1mm,
  fontupper=\small,
}

% 答题区 — 无边框，纯留白
\newcommand{\answerarea}[1][40mm]{%
  \vspace{2mm}%
  \begin{tcolorbox}[
    enhanced, breakable,
    colback=white,
    colframe=white,
    boxrule=0pt,
    height=#1,
    valign=top,
    bottom=0mm,
    after skip=0mm,
  ]
  \end{tcolorbox}%
}

% ---------- 字体设置 ----------
% exam-zh (ctexbook) already sets CJK fonts via ctex.
% Only override if specific fonts are needed.
% macOS: Songti SC, STHeiti, STFangsong
% Linux: SimSun, SimHei, FangSong

% ---------- 页面设置 ----------
% exam-zh (ctexbook) already loads geometry.
% Override margins if needed:
% \geometry{top=16mm, bottom=16mm, left=14mm, right=14mm}

\examsetup{
  question/show-answer = true,
  fillin/show-answer = true,
  solution/show-solution = show-stay,
}

% ---------- 教师版解答样式 ----------
\newtcolorbox{solutionblock}{
  enhanced, breakable,
  colback=edu-blue-bg,
  colframe=edu-blue!25,
  boxrule=0pt,
  borderline west={2pt}{0pt}{edu-blue!50},
  arc=0pt,
  left=3mm, right=3mm, top=2mm, bottom=2mm,
  before skip=2mm,
  after skip=3mm,
  fontupper=\small,
}

% 解答步骤编号
\newcounter{solstep}
\newcommand{\solstep}[2]{%
  \stepcounter{solstep}%
  \par\medskip
  \noindent\hangindent=6mm
  {\color{edu-blue}\bfseries\textcircled{\scriptsize\thesolstep}\enspace #1}\enspace #2\par
}

\begin{document}

\section*{点到直线距离与“线圆位置关系”的互推妙用 · 专题练习（教师版）}

\section*{一、选择题（每题 4 分，共 16 分）}


\begin{question}[points=4]
  已知 $\odot O$ 的半径为 $5\text{cm}$，圆心 $O$ 到直线 $l$ 的距离 $d = 4\text{cm}$，则直线 $l$ 与 $\odot O$ 的位置关系是
  \begin{choices}
    \item 相交
    \item 相切
    \item 相离
    \item 无法确定
  \end{choices}
  \paren[A]
  \begin{solutionblock}
    因为圆心 $O$ 到直线 $l$ 的距离 $d = 4\text{cm}$，半径 $r = 5\text{cm}$。因为 $d < r$，根据直线与圆的位置判定定理，直线 $l$ 与 $\odot O$ 相交。
  \end{solutionblock}
\end{question}



\begin{question}[points=4]
  若直线 $l$ 与半径为 $r$ 的 $\odot O$ 相切，则直线 $l$ 与 $\odot O$ 的公共点个数为
  \begin{choices}
    \item $0$ 个
    \item $1$ 个
    \item $2$ 个
    \item 无数个
  \end{choices}
  \paren[B]
  \begin{solutionblock}
    直线与圆相切的几何定义为直线与圆只有一个公共点。
  \end{solutionblock}
\end{question}



\begin{question}[points=4]
  在 $Rt\triangle ABC$ 中，$\angle C = 90^\circ$，$AC = 6$，$BC = 8$。以点 $C$ 为圆心，以 $5$ 为半径画 $\odot C$，则斜边 $AB$ 与 $\odot C$ 的位置关系是
  \begin{choices}
    \item 相交
    \item 相切
    \item 相离
    \item 无法确定
  \end{choices}
  \paren[A]
  \begin{solutionblock}
    在 $Rt\triangle ABC$ 中，根据勾股定理斜边 $AB = \sqrt{6^2 + 8^2} = 10$。利用等面积法求圆心 $C$ 到斜边 $AB$ 的垂线段距离 $d = \frac{AC \cdot BC}{AB} = \frac{6 \times 8}{10} = 4.8$。因为距离 $d = 4.8 < r = 5$（即 $d < r$），所以斜边 $AB$ 与 $\odot C$ 相交。
  \end{solutionblock}
\end{question}



\begin{question}[points=4]
  已知 $\odot O$ 的半径为 $3$，圆心 $O$ 到直线 $l$ 的距离为 $d$。若直线 $l$ 与 $\odot O$ 相交，则 $d$ 的取值范围是
  \begin{choices}
    \item $d < 3$
    \item $0 \le d < 3$
    \item $d > 3$
    \item $0 < d \le 3$
  \end{choices}
  \paren[B]
  \begin{solutionblock}
    因为直线与圆相交，所以 $d < r = 3$。又因为圆心到直线的距离 $d$ 作为线段长度，必然是非负数，即 $d \ge 0$。综上，距离 $d$ 的取值范围是 $0 \le d < 3$。
  \end{solutionblock}
\end{question}


\section*{二、填空题（每题 4 分，共 8 分）}


\begin{question}[points=4]
  在等腰 $\triangle ABC$ 中，$AB = AC = 5$，$BC = 6$。若以点 $A$ 为圆心，以 $r$ 为半径的 $\odot A$ 与底边 $BC$ 相切，则半径 $r = $ \underline{\quad\quad}。
  \paren[$4$]
  \begin{solutionblock}
    作 $AD \perp BC$ 于点 $D$。因为 $AB = AC = 5$，$BC = 6$，根据等腰三角形三线合一特性，$D$ 为 $BC$ 的中点，所以 $BD = 3$。在 $Rt\triangle ABD$ 中，根据勾股定理得高 $AD = \sqrt{AB^2 - BD^2} = \sqrt{5^2 - 3^2} = 4$。因此圆心 $A$ 到直线 $BC$ 的距离 $d = AD = 4$。因为 $\odot A$ 与底边 $BC$ 相切，所以 $r = d = 4$。
  \end{solutionblock}
\end{question}



\begin{question}[points=4]
  在 $Rt\triangle ABC$ 中，$\angle C = 90^\circ$，$AC = 3$，$BC = 4$。以点 $C$ 为圆心，以 $r$ 为半径画 $\odot C$。若斜边 $AB$ 与 $\odot C$ 相离，则半径 $r$ 的取值范围是 \underline{\quad\quad}。
  \paren[$0 < r < 2.4$]
  \begin{solutionblock}
    在 $Rt\triangle ABC$ 中，根据勾股定理斜边 $AB = \sqrt{3^2 + 4^2} = 5$。由面积法得圆心 $C$ 到斜边 $AB$ 的垂直距离 $d = \frac{3 \times 4}{5} = 2.4$。因为斜边 $AB$ 与 $\odot C$ 相离，所以距离大于半径，即 $r < d = 2.4$。又因为圆的半径必须为正数（$r > 0$），所以半径 $r$ 的取值范围是 $0 < r < 2.4$。
  \end{solutionblock}
\end{question}


\section*{三、解答题（每题 10 分，共 20 分）}

\needspace{8\baselineskip}

\begin{problem}[points=10]
  在 $\triangle ABC$ 中，$\angle C = 90^\circ$，$AC = 5$，$BC = 12$。以点 $C$ 为圆心，作以 $r$ 为半径的 $\odot C$。
\begin{enumerate}[label=(\arabic*)]
  \item 若 $\odot C$ 与斜边 $AB$ 相切，求半径 $r$ 的值；
  \item 若半径 $r = 5.2$，判断斜边 $AB$ 与 $\odot C$ 的位置关系。
\end{enumerate}

  \setcounter{solstep}{0}
  \begin{solutionblock}
    \textbf{答案：}(1) 半径 $r = \frac{60}{13}$ ； (2) 直线 $AB$ 与 $\odot C$ 相交。\par\medskip
    \solstep{求斜边}{在 $Rt\triangle ABC$ 中，由勾股定理：$AB = \sqrt{AC^2 + BC^2} = \sqrt{5^2 + 12^2} = 13$}
    \solstep{求高（距离 d）}{利用等面积法 $AC \cdot BC = AB \cdot d \implies d = \frac{5 \times 12}{13} = \frac{60}{13}$}
    \solstep{解答第 (1) 问}{因为直线与圆相切，所以 $r = d = \frac{60}{13}$。}
    \solstep{解答第 (2) 问}{圆心到直线的实际距离 $d = \frac{60}{13} \approx 4.62$。因为 $d = 4.62 < 5.2 = r$，即 $d < r$，根据判定定理，斜边 $AB$ 与 $\odot C$ 相交。}
  \end{solutionblock}
\end{problem}


\needspace{8\baselineskip}

\begin{problem}[points=10]
  在平面直角坐标系中，以点 $A(3, 4)$ 为圆心，作以 $r$ 为半径的 $\odot A$。
\begin{enumerate}[label=(\arabic*)]
  \item 若 $\odot A$ 与 $x$ 轴相切，求半径 $r$ 的值；
  \item 若 $\odot A$ 与 $y$ 轴相交，求半径 $r$ 的取值范围。
\end{enumerate}

  \setcounter{solstep}{0}
  \begin{solutionblock}
    \textbf{答案：}(1) 半径 $r = 4$ ； (2) 半径 $r > 3$。\par\medskip
    \solstep{确定圆心到两轴的距离}{圆心为 $A(3, 4)$。则圆心到 $x$ 轴的距离 $d_1 = |y_A| = 4$；圆心到 $y$ 轴的距离 $d_2 = |x_A| = 3$。}
    \solstep{解答第 (1) 问}{因为 $\odot A$ 与 $x$ 轴相切，所以 $r = d_1 = 4$。}
    \solstep{解答第 (2) 问}{因为 $\odot A$ 与 $y$ 轴相交，根据判定定理有 $d_2 < r \implies 3 < r$，即半径 $r$ 的取值范围是 $r > 3$。}
  \end{solutionblock}
\end{problem}


\clearpage
\section*{参考答案}


\begin{keyideabox}
  \textbf{选择题}
  1. A（因为 $d = 4 < r = 5$，相交）
2. B（相切时有且仅有 1 个公共点）
3. A（$AB=10, d=4.8 < r=5$，相交）
4. B（$d < 3$ 且距离 $d \ge 0$，因此 $0 \le d < 3$）

\end{keyideabox}



\begin{keyideabox}
  \textbf{填空题}
  1. $4$（高线 $AD = \sqrt{5^2 - 3^2} = 4 = r$）
2. $0 < r < 2.4$（斜边高 $d = 2.4 > r$ 且半径 $r > 0$）

\end{keyideabox}



\begin{keyideabox}
  \textbf{解答题}
  第 1 题：(1) 斜边 $AB = 13$，高 $d = \frac{60}{13}$，因为相切，半径 $r = d = \frac{60}{13}$；
(2) 当 $r = 5.2$ 时，因为 $d = \frac{60}{13} \approx 4.62 < 5.2$，所以斜边 $AB$ 与 $\odot C$ 相交。
第 2 题：(1) 到 $x$ 轴距离 $d_1 = 4$，相切时 $r = 4$；
(2) 到 $y$ 轴距离 $d_2 = 3$，相交时 $d_2 < r \implies r > 3$。

\end{keyideabox}



\end{document}
